% sec-06: the quantitative case.  Figures 16 (d2), 17 (graphviz).
% A funder buys a trial, not a plan.  This section is what they are buying.
\section{The Clinical Evidence a Funder Is Buying}

\subsection{Six Quantities a Reviewer Can Check}

\draftinstr{One short paragraph stating the discipline before any number: every
quantity below is published or deposited, every one carries the limitation its
own authors stated, and no in-silico result is presented as a clinical one.
Then set Table 15 with six rows and six columns: quantity, test arm, comparator,
evidence tier, date, and the limitation in the source authors' own words. Every
row is in
\appfile{funding/pdac-funding-applications/final-apply/sections/sec-05-trial-evidence.tex}
and must be carried unchanged. Cite \texttt{rasolute302},
\texttt{qspmetpancre}, \texttt{daraxsims}, \texttt{daraxident},
\texttt{asmevv40} and \texttt{ichm15}.}

\figslot{d2-type / grid with intervals}{7.8}{Six quantities a reviewer can check
against a published source,\\each beside its comparator and its stated
limitation. Two are in\\silico, one is a digital twin, and three are trial or
registry.}

\draftinstr{After Figure 16, the chronology paragraph and the caveat that must
travel with it. The 2025 simulated ratio was 2.4-fold and the 2026 observed
ratio was 2.0-fold. State the proximity, then state in the same paragraph that
it is a chronology observation and a hypothesis-supporting one and not a
validation claim, and name the three material differences: 1000 simulated
against 241 enrolled, a combination against a single agent, and KRAS G12C
against a primarily G12D and G12V population.}

\subsection{The Trial Being Funded}

\draftinstr{One paragraph of trial design, from
\appfile{trial-protocol/} and \appfile{trial-ind/}: open-label, single-arm, 3+3
escalation, up to 18 treated participants, perioperative daraxonrasib with an
eight-arm robotic pancreaticoduodenectomy and an on-premises model confined to
advisory output. Carry the device numbers a reviewer will ask for: 56 degrees of
freedom, 640 sensor channels, 3 N per-arm and 18 N cumulative force caps, 3 ms
cross-arm arm stop and 500 ms system stop, and the five-vessel no-fly gate.
Include the positioning correction: daraxonrasib is investigational and already
in Phase 3 evaluation and is nowhere described as first in human; the
supportable novelty claim concerns the integrated surgical and advisory
workflow. Source: \appfile{funding/potential-partners/UC-San-Diego/}.}

\subsection{Where Each Number Lands}

\figslot{graphviz-type / record chain}{8.0}{Five published quantities and the
exact filing section each one\\reaches. Every edge carries its link type, and
every chain ends\\at a milestone whose artifact a program officer can request.}

\draftinstr{After Figure 17, one paragraph on the five link types and why each
edge is labelled: naming the link type is what stops the chain being read as an
inheritance. A tier 1 result does not become a tier 3 result by travelling along
an arrow. Source:
\appfile{funding/capitalization-plan/graphviz/fig-17-evidence-chain-records.gv.md}.}

\subsection{What the Work Cost, and What It Costs Elsewhere}

\draftinstr{Set Table 16: three work products against their industry cost and
schedule benchmarks and against the author's own estimated figures, carried
unchanged from
\appfile{funding/pdac-funding-applications/final-apply/sections/sec-05-trial-evidence.tex}.
Label the author figures as estimates from the source set and not as audited
accounts, in the caption, in the same words the parent work uses. Close with one
sentence connecting this table to Figure 3: the same cost discipline that
produced a \$36,330 virtual trial produces a 7.5 percent indirect rate.}
